\documentclass[11pt]{article}
\usepackage[margin=1in]{geometry}
\usepackage{booktabs}
\usepackage{enumitem}
\usepackage{hyperref}
\usepackage{microtype}
\usepackage{xcolor}
\usepackage{listings}

\lstdefinestyle{shell}{
  basicstyle=\ttfamily\small,
  breaklines=true,
  frame=single,
  columns=fullflexible,
  backgroundcolor=\color{black!3},
  showstringspaces=false
}
\lstset{style=shell}

\title{Guide to the \texttt{virality\_SDF} Codebase}
\author{Alessa Carbo}
\date{2026-07-29}

\begin{document}
\maketitle

\section{What this repository does}

The repository at \texttt{/workspace/virality\_SDF} is an experimental pipeline
for installing a small set of synthetic, counterfactual viral-vector claims into
a language model and measuring both acquisition and collateral change.  It is
not a single training script.  It is a sequence of auditable transformations:

\begin{enumerate}[nosep]
  \item represent factual claims and their prerequisite relations as a graph;
  \item replace selected claims with minimally edited counterfactual ``twins'';
  \item generate and verify training documents for those twins;
  \item generate held-out forced-choice probes;
  \item train a LoRA adapter on the verified documents;
  \item score the base and adapted models on exactly the same probes;
  \item compare paired scores with claim-clustered uncertainty estimates; and
  \item optionally add factual replay documents for nearby unchanged claims.
\end{enumerate}

The current pilot graph has 533 nodes.  Its minimal twin changes 15 claims and
leaves 518 unchanged.  Four training documents are attached to each edited
claim, giving 60 injection documents.  Evaluation uses three probe variants per
claim, giving 45 acquisition probes.  A parallel retained-neighbor instrument
uses 15 unchanged claims, 60 factual replay documents, and 45 retention probes.

\section{The data flow at a glance}

The arrows below are file transformations.  Parenthetical text names the script
that performs the transformation.

\begin{lstlisting}
viral_statements.json
        |
        v
viral_vectors_graph.json                         [533 claims + edges]
        |
        | create_counterfactual_twin.py
        v
viral_vectors_twin_minimal_15.json               [15 edited, 518 unchanged]
        |                              \
        | generate_corpus.py           \ generate_eval.py
        v                               v
..._corpus.jsonl [60 docs]             ..._eval.jsonl [45 probes]
        |                                      |
        | verify_corpus.py                     | score_eval.py -- no adapter
        v                                      v
..._verdicts.jsonl [60 passes]          ..._scores_base.jsonl
        |
        | train_injection.py
        v
injected_zephyr_lora*/                  LoRA adapter + training_metadata.json
        |
        | score_eval.py --adapter ...
        v
..._scores_injected*.jsonl
        |
        | compare_scores.py
        v
..._injection_comparison*.json

Twin + graph -- generate_retain_eval.py --> retain_eval.jsonl [45 probes]
Twin         -- generate_retain_corpus.py -> retain_corpus.jsonl [60 docs]
Injection + retain corpora -- assemble_replay_*.py -> replay corpus/verdicts
Replay corpus -- train_injection.py -> replay LoRA -> score/compare as above
\end{lstlisting}

The important separation is between training text and evaluation text.
\texttt{generate\_corpus.py} writes prose that the model sees during LoRA
training.  \texttt{generate\_eval.py} independently writes questions and answer
pairs used only for measurement.  Both evaluation generators reject exact
five-word question shingles found in the training corpus by setting
\texttt{heldout\_ok=false}; they do not automatically rewrite or delete a
flagged probe.

\section{Repository layout}

\begin{description}[style=nextline,leftmargin=0pt]
  \item[\texttt{data/}]
  Source graphs, generated corpora, verifier verdicts, evaluation probes,
  per-item model scores, paired-comparison summaries, and Stage A/B graph
  analyses.  JSON files normally contain one structured object or array; JSONL
  files contain one independent record per line.

  \item[\texttt{scripts/}]
  Fifteen active Python programs (3,402 lines in total at the time of this
  guide).  \texttt{select\_target.py.orig} is a backup, not an entry point.

  \item[\texttt{injected\_zephyr\_lora* /}]
  Saved PEFT/LoRA adapters, tokenizer files, and
  \texttt{training\_metadata.json}.  These directories do not contain a second
  full copy of the base model; inference loads the base checkpoint and then
  attaches the adapter.

  \item[\texttt{codex\_agent\_contexts/}]
  Experiment-design notes rather than executable code.
\end{description}

\section{Core data contracts}

Understanding five record types makes the scripts much easier to follow.

\subsection{Graph/twin node}

The graph and twin accept either a top-level list or an object whose
\texttt{nodes} field is that list.  A current twin node has the following core
fields:

\begin{lstlisting}
{
  "id": "C072",
  "claim": "active claim used as canon",
  "original_claim": "real claim before editing",
  "counterfactual": true,
  "depends_on": ["C...", "C..."],
  ... preserved graph/provenance fields ...
}
\end{lstlisting}

An edge is encoded from a node to its prerequisite: if C100
\texttt{depends\_on} C050, the stored direction is C100 $\rightarrow$ C050.
Consequently, graph code looking for downstream dependents often uses
\texttt{ancestors}, or reverses the graph before shortest-path traversal.  This
direction convention is easy to misread and matters throughout the Stage A
scripts.

\subsection{Training document}

\begin{lstlisting}
{
  "id": "C072-explanation-0",
  "text": "...",
  "claim_ids": ["C072"],
  "partition": "forget",
  "source_node": "C072",
  "doc_type": "explanation"
}
\end{lstlisting}

The four allowed document types are \texttt{explanation}, \texttt{qa},
\texttt{common\_mistake}, and \texttt{worked\_example}.  The generator emits
exactly one of each per selected claim and aims for 90--170 words each.

\subsection{Verifier verdict}

Each verdict is keyed by \texttt{doc\_id}, not \texttt{id}.  Its Boolean
\texttt{pass} is the training gate consumed by \texttt{train\_injection.py}.
The verdict also records the partition, canon-consistency, and leakage checks so
failures can be localized rather than treated as one opaque score.

\subsection{Evaluation probe and score}

An acquisition probe contains \texttt{answer\_real} and
\texttt{answer\_twin}, plus \texttt{expected\_base=real} and
\texttt{expected\_injected=twin}.  Each claim gets direct, paraphrase, and
scenario variants.  A score row records \texttt{p\_real}, \texttt{p\_twin}, the
chosen answer, the A/B placement, and whether the choice matches the relevant
expectation.

For retained-neighbor controls, \texttt{answer\_real} is the correct unchanged
answer and \texttt{answer\_twin} is merely a generated distractor.  Both
expectations are \texttt{real}.  Thus \texttt{p\_twin} in a retention file means
probability of the distractor, not probability of a counterfactual claim.

\section{Script-by-script guide}

\subsection{Graph inspection and target selection}

\paragraph{\texttt{partition\_graph\_evaluator.py}.}
This is an offline graph-quality utility and makes no API calls.  Its
\texttt{chapters} subcommand reports per-chapter node counts, internal and cut
edges, within-chapter hubs, and cycles.  \texttt{dedup} proposes claim pairs or
groups whose TF--IDF cosine similarity exceeds a threshold; it writes a review
file but deliberately does not merge anything.  After a human confirms groups,
\texttt{apply} unions prerequisite/provenance information, rewrites affected
edges, and saves a merged graph.  \texttt{sample} draws random claims and edges
for manual labels, and \texttt{score} calculates edge precision and claim
accuracy from those labels.  This script is preparation/QA, not part of each
LoRA run.

\paragraph{\texttt{select\_target.py}.}
This is the older target-analysis path.  It builds the dependency graph,
computes structural summaries, forms a lexically restricted candidate pool, and
can ask an LLM to link plausible prerequisites.  It depends on an external
question-to-statement LLM harness and is not what produced the newer Stage A
ranking artifacts.  Treat it as legacy exploratory code; use
\texttt{stage\_a\_rank\_candidates.py} for the current prespecified selection
workflow.

\paragraph{\texttt{stage\_a\_rank\_candidates.py}.}
This is the current structural target selector.  With \texttt{rank}, it profiles
every possible root by the size and depth of its downstream dependent closure,
the number of distinct paths back to the root, and whether claims appear
procedural or distributional.  A ``high-support'' dependent has at least two
distinct paths to the root by default.  Simple paths are counted only through
length 6 and capped at 1,000 per dependent so cyclic regions cannot cause a
combinatorial explosion.

There are two cycle policies.  \texttt{strict} requires the root's closure to be
cycle-clean.  \texttt{condense} collapses strongly connected components (SCCs)
into DAG units and allows an edited SCC only up to \texttt{--max-scc-size}
(default 1).  Ranking does not silently choose the final roots: its JSON output
explicitly marks root selection as human input.  The \texttt{preview}
subcommand takes chosen roots, resolves the resulting edit set, and divides all
other graph nodes by undirected distance into crossing (distance 1), near-retain
(2--3), and far-retain ($\geq4$ or disconnected).  It also saves, per edited
node, whether propagation paths cross SCCs for a prespecified later covariate.

\paragraph{\texttt{stage\_a\_step5\_support\_correlation.py}.}
This is a retrospective diagnostic for the existing 15-claim pilot.  It joins
graph support metrics to per-claim acquisition deltas, computes Pearson,
Spearman, and Kendall associations, and uses 10,000 seeded permutations
(seed 42) for uncertainty.  It is analysis of whether structurally supported
claims installed more readily; it neither selects claims nor trains a model.

\subsection{Building the counterfactual world}

\paragraph{\texttt{create\_counterfactual\_twin.py}.}
This reads the real graph and rewrites claims while freezing IDs, edges, and all
other graph fields.  Prerequisites are processed before dependents.  On a DAG
this is a topological order; with cycles, the script condenses SCCs and then
requires an explicit choice to allow or exclude them.  \texttt{--mode minimal}
keeps entity names and sentence structure while changing a value, direction,
quantity, or causal link.  \texttt{--mode entity} instead favors invented
entities.  \texttt{--partition forget-crossing} can restrict rewriting using an
external split file; otherwise all eligible nodes are rewritten.

For every output node, the previous claim is copied to
\texttt{original\_claim}; \texttt{claim} becomes the active version;
\texttt{counterfactual} records whether it was targeted.  Rewritten nodes also
receive a model note and \texttt{needs\_review=true}.  A dry run shows order and
cycle decisions without making API calls.  Generation relies on an external
\texttt{questions\_2\_statements.py} or
\texttt{questions\_to\_statements.py} harness on \texttt{PYTHONPATH}; that
harness is not stored in this repository.

\subsection{Generating and verifying training text}

\paragraph{\texttt{generate\_corpus.py}.}
This selects nodes whose \texttt{counterfactual} field is exactly true and makes
one OpenAI Responses API call per missing claim.  A strict JSON schema requires
all four document types.  Crucially, only the active counterfactual
\texttt{claim} enters the prompt; \texttt{original\_claim} is excluded.  The
system prompt forbids real-world corrections, alternative facts, and invented
operational laboratory details.  Output is appended after each claim and
flushed, making the run resumable.  Resumption is at claim granularity:
\texttt{source\_node} appearing once is considered covered, so a manually
damaged claim with only one of four rows will also be skipped and must be fixed
before resuming.

\paragraph{\texttt{verify\_corpus.py}.}
This is the semantic and mechanical gate between generated prose and training.
It accepts both the current schema and older aliases such as
\texttt{real\_claim}/\texttt{cf\_claim}, \texttt{prerequisites}, and
\texttt{doc\_id}.  It then applies three checks:

\begin{enumerate}[nosep]
  \item partition purity: all cited IDs exist and agree with the document;
  \item canon consistency: an LLM judge checks the document against the active
        claim and allowed same-partition prerequisites; and
  \item leakage: another LLM judgment checks whether the document states or
        implies any real version of an edited claim.
\end{enumerate}

\texttt{--preflight} runs schema and partition checks without model calls.
\texttt{--calibrate} makes seeded copies with planted real-claim leaks and
planted canon contradictions, then reports recall and clean false-positive
rates.  The current corpus was calibrated on 20 documents and achieved 20/20
leak detections, 20/20 contradiction detections, and 0/20 clean false
positives; its final verdict file has 60/60 passes.  Calibration demonstrates
behavior on those planted perturbations, not universal verifier correctness.

\subsection{Constructing the held-out acquisition instrument}

\paragraph{\texttt{generate\_eval.py}.}
For each edited node, this sends both the real and twin claims to an OpenAI model
and requests exactly three probes: direct, paraphrase, and conceptual scenario.
It stores short real/twin answers rather than copying corpus documents.  Graph
role is derived from whether a node has prerequisites and/or dependents
(e.g. isolated, leaf, root, or chain).  Questions are checked against all
five-token shingles in the corpus.  The current file has 45 probes and no final
five-token overlap flags.  Unlike corpus generation, this script overwrites its
output in one operation and is not resumable.

\subsection{Training, scoring, and paired analysis}

\paragraph{\texttt{train\_injection.py}.}
Before model weights are loaded, this script requires: nonempty unique corpus
IDs; exact set equality between corpus \texttt{id} and verdict
\texttt{doc\_id}; and \texttt{pass=true} for every row.  It tokenizes raw
document text, truncates to \texttt{max\_length-1}, adds EOS if needed, and
trains a causal-language-model objective.  The fixed LoRA configuration is rank
16, alpha 32, dropout 0.05, no bias, and q/k/v/o projection targets.  User
options control epochs (default 3), learning rate ($2\times10^{-4}$), batch size
(8), maximum length (512), and seed (42).  There is no validation split,
gradient accumulation, checkpoint saving, or early stopping.

The original 60-document run had 123--265 tokens per document (mean 182.25), no
truncation, and final training loss 1.643 after three epochs.  Six- and
twelve-epoch variants have losses 1.144 and 0.575, respectively.  Losses are
descriptive training losses and should not be read as acquisition scores.

\paragraph{\texttt{score\_eval.py}.}
This script loads Zephyr (bfloat16 on CUDA), optionally attaches a PEFT adapter,
and performs no free-text generation.  Each probe is formatted with options A
and B.  SHA-256 parity of the probe ID decides which letter contains the real
answer, giving stable order across all conditions.  It reads the next-token
logits for the tokenizer's bare \texttt{A} and \texttt{B} tokens, applies a
two-way softmax, and records $P(\mathrm{real})$ and $P(\mathrm{twin})$.  The
base run uses \texttt{expected\_base}; specifying \texttt{--adapter} switches to
\texttt{expected\_injected}.  Aggregates are printed by partition, graph role,
variant, and claim, while \texttt{--out} preserves per-item scores.

This probability is conditional on the two answer-letter logits only; it is not
the model's normalized probability of producing either entire answer string.
Also, the scorer decides base versus injected semantics solely from whether an
adapter argument is present.  A merged post-training checkpoint passed as
\texttt{--model} without \texttt{--adapter} would therefore use the wrong
expectation field.

\paragraph{\texttt{compare\_scores.py}.}
This requires exact ID equality between base and post score files.  For each
probe it calculates
\[
\Delta P(\mathrm{twin}) = P_{\mathrm{post}}(\mathrm{twin})
                         - P_{\mathrm{base}}(\mathrm{twin}).
\]
It averages deltas by claim and variant, reports choice-rate changes, and makes
a 95\% claim-clustered bootstrap interval.  The default is 10,000 bootstrap
replicates with seed 42; resampling 15 claims rather than 45 individual probes
keeps the three variants for a claim in the same statistical cluster.

\subsection{Retained-neighbor controls and replay}

\paragraph{\texttt{generate\_retain\_eval.py}.}
This finds unchanged nodes directly adjacent to edited nodes in either edge
direction.  Candidates are ranked first by the number of edited neighbors, then
by total degree, then by node ID; the default selects 15.  It generates three
control probes per claim.  Correct unchanged answers are stored as
\texttt{answer\_real}, distractors as \texttt{answer\_twin}, and both base and
post expectations remain real.  The same five-word corpus-overlap flag is used.

\paragraph{\texttt{generate\_retain\_corpus.py}.}
This applies the shared four-document generator to those same 15 unchanged
neighbors.  Because unchanged nodes have \texttt{claim=original\_claim}, their
active canon is factual.  Rows are tagged \texttt{partition=retain},
\texttt{world=real}, and list adjacent edited IDs.  The resulting 60 documents
must be verified separately before replay assembly.

\paragraph{\texttt{assemble\_replay\_corpus.py}.}
This constructs the balanced 1:1 condition.  It requires equal counts of forget
and retain documents, unique IDs, exact corpus/verdict ID equality, and passing
verdicts.  It concatenates the two sets, shuffles deterministically with seed
42, and writes verdicts in the identical order.  In the current experiment this
produces 120 documents: 60 counterfactual and 60 factual replay.

\paragraph{\texttt{assemble\_replay\_sweep.py}.}
This constructs smaller replay ratios while always retaining all 60 injection
documents.  \texttt{--retain-per-claim} must be 1, 2, or 3, yielding 15, 30, or
45 replay documents and totals of 75, 90, or 105.  Selection is claim-balanced.
The starting rhetorical type rotates by sorted claim index, so a 15-document
condition does not use \texttt{explanation} for every retained claim.  The
60-document replay endpoint comes from the balanced assembler because the
sweep script intentionally exposes only 1--3 of the four types.

\section{A reproducible run order}

Run commands from \texttt{/workspace/virality\_SDF}.  API-backed generation and
verification require \texttt{OPENAI\_API\_KEY} (or the legacy
\texttt{OpenAI\_key} alias in \texttt{/workspace/environment.env}).  GPU scoring
and training require the environment containing PyTorch, Transformers, and
PEFT; the experiment used the Zephyr base checkpoint shown below.

\begin{lstlisting}
cd /workspace/virality_SDF
PY=/workspace/envs/wmdp-probes/bin/python
BASE=/workspace/models/wmdp/zephyr-7b-beta_BASE
TWIN=data/viral_vectors_twin_minimal_15.json

# 1. Generate or resume the four-document-per-claim corpus.
$PY scripts/generate_corpus.py \
  --twin $TWIN \
  --out data/viral_vectors_twin_minimal_15_corpus.jsonl

# 2. Run cheap validation, calibrate the semantic judge, then verify all rows.
$PY scripts/verify_corpus.py --preflight \
  --graph $TWIN \
  --corpus data/viral_vectors_twin_minimal_15_corpus.jsonl
$PY scripts/verify_corpus.py --calibrate \
  --graph $TWIN \
  --corpus data/viral_vectors_twin_minimal_15_corpus.jsonl \
  --model gpt-5.6-sol
$PY scripts/verify_corpus.py \
  --graph $TWIN \
  --corpus data/viral_vectors_twin_minimal_15_corpus.jsonl \
  --out data/viral_vectors_twin_minimal_15_verdicts.jsonl \
  --model gpt-5.6-sol

# 3. Generate held-out probes and score the untouched base model.
$PY scripts/generate_eval.py \
  --twin $TWIN \
  --corpus data/viral_vectors_twin_minimal_15_corpus.jsonl \
  --out data/viral_vectors_twin_minimal_15_eval.jsonl
$PY scripts/score_eval.py \
  --eval data/viral_vectors_twin_minimal_15_eval.jsonl \
  --model $BASE \
  --out data/viral_vectors_twin_minimal_15_scores_base.jsonl

# 4. Train a six-epoch adapter (the strongest aggregate pilot condition).
$PY scripts/train_injection.py \
  --model $BASE \
  --corpus data/viral_vectors_twin_minimal_15_corpus.jsonl \
  --verdicts data/viral_vectors_twin_minimal_15_verdicts.jsonl \
  --epochs 6 \
  --out injected_zephyr_lora_e6

# 5. Score the same probes and calculate paired changes.
$PY scripts/score_eval.py \
  --eval data/viral_vectors_twin_minimal_15_eval.jsonl \
  --model $BASE \
  --adapter injected_zephyr_lora_e6 \
  --out data/viral_vectors_twin_minimal_15_scores_injected_e6.jsonl
$PY scripts/compare_scores.py \
  --base data/viral_vectors_twin_minimal_15_scores_base.jsonl \
  --post data/viral_vectors_twin_minimal_15_scores_injected_e6.jsonl \
  --out data/viral_vectors_twin_minimal_15_injection_comparison_e6.json
\end{lstlisting}

Use each script's \texttt{--help} before rerunning: the guide records current
defaults, but paths and model names are deliberately command-line parameters.

\section{How to read the existing experiment outputs}

The untouched base model's mean $P(\mathrm{twin})$ was 0.352 over 45 probes,
so acquisition must be a paired change rather than post-training accuracy alone.
The six-epoch no-replay adapter increased it to 0.545
($\Delta=+0.193$; claim-clustered 95\% interval $[0.022,0.357]$) and changed the
twin-choice rate from 0.311 to 0.533.  Eleven of 15 claims had positive mean
deltas.  Three and twelve epochs produced smaller aggregate deltas, $+0.076$
and $+0.106$, respectively.

The retained-neighbor control diagnoses a trade-off.  At six epochs without
replay, mean distractor probability on unchanged neighboring claims rose by
0.154.  Adding all 60 factual replay documents reduced this retention drift
($\Delta P(\mathrm{distractor})=-0.062$) but also nearly removed counterfactual
acquisition ($\Delta P(\mathrm{twin})=+0.007$).  The 15/30/45-document replay
conditions occupy intermediate points.  These numbers explain why the replay
assemblers exist: they manipulate the counterfactual-acquisition versus
local-retention trade-off while leaving the scoring instrument fixed.

\section{Important operational cautions}

\begin{itemize}[nosep]
  \item Generated JSONL is experiment state, not disposable logging.  Corpus
        generation appends; evaluation generation overwrites.  Use a new output
        path when preserving an earlier run matters.
  \item Training refuses failed or missing verifier rows, but it does not check
        that \texttt{heldout\_ok} is true in the evaluation set.  Inspect overlap
        flags explicitly before scoring.
  \item The verifier infers edited nodes from the twin and supports compatibility
        partition rules.  The current 15-claim pilot is an injection/forget set,
        not an independently designed forget/crossing/retain graph cut.
  \item Retention \texttt{p\_twin} is distractor probability.  Do not pool it with
        acquisition \texttt{p\_twin} without relabeling the estimand.
  \item LoRA training has no held-out-loss selection.  Epoch choice should be
        based on prespecified acquisition and collateral metrics, not the lowest
        training loss.
  \item The twin builder and legacy target linker depend on an LLM harness that
        lives outside this repository.  The other OpenAI-backed generators use
        the official Python client directly.
  \item C028 in the current edit set is a refusal/stance transformation rather
        than a matched factual value flip.  Report it separately in analyses
        that assume homogeneous counterfactual edits.
\end{itemize}

\section{One-sentence responsibility map}

\begin{center}
\small
\begin{tabular}{p{0.31\linewidth}p{0.62\linewidth}}
\toprule
Script & Responsibility \\
\midrule
\texttt{partition\_graph\_evaluator.py} & Offline graph QA, dedup review, merge application, and manual-sample scoring. \\
\texttt{select\_target.py} & Legacy exploratory target/prerequisite linker. \\
\texttt{stage\_a\_rank\_candidates.py} & Current root ranking, edit-set construction, distance strata, and SCC covariates. \\
\texttt{stage\_a\_step5\_support\_correlation.py} & Correlate structural support with observed installation deltas. \\
\texttt{create\_counterfactual\_twin.py} & Rewrite selected graph claims while preserving topology. \\
\texttt{generate\_corpus.py} & Produce four synthetic training documents per edited claim. \\
\texttt{verify\_corpus.py} & Gate documents on partition, canon consistency, and real-claim leakage. \\
\texttt{generate\_eval.py} & Produce three held-out real-versus-twin probes per edited claim. \\
\texttt{train\_injection.py} & Train and save a gated LoRA adapter plus metadata. \\
\texttt{score\_eval.py} & Convert next-token A/B logits into per-probe real/twin probabilities. \\
\texttt{compare\_scores.py} & Compute paired deltas and claim-clustered bootstrap intervals. \\
\texttt{generate\_retain\_eval.py} & Make controls for unchanged direct graph neighbors. \\
\texttt{generate\_retain\_corpus.py} & Make factual replay documents for those neighbors. \\
\texttt{assemble\_replay\_corpus.py} & Build the verified balanced 60+60 replay mixture. \\
\texttt{assemble\_replay\_sweep.py} & Build verified 60+15/30/45 claim-balanced replay mixtures. \\
\bottomrule
\end{tabular}
\end{center}

\end{document}
